% ==================================================================
%  WEEK 8
% ==================================================================
\chapter{Containers, Docker Virtualization, and Sensor-Cloud Infrastructure}

\section{Introduction and Week Overview}
Week 8 transitions from heavy hypervisor-based hardware virtualization into lightweight operating-system-level virtualization (Containers and Docker), followed by the architecture of Sensor-Cloud computing. The key modules examined include:
\begin{enumerate}[leftmargin=*]
  \item \textbf{Containers vs Virtual Machines:} Kernel-sharing mechanics, isolation primitives, architectural tradeoffs, and performance benchmarking.
  \item \textbf{Docker Ecosystem and Architecture:} The Docker engine, copy-on-write image layering (UnionFS/OverlayFS), Dockerfile build directives, and container lifecycle CLI commands.
  \item \textbf{Container Orchestration Primitives:} High-level role of Kubernetes and distributed configuration state management via \texttt{etcd}.
  \item \textbf{Sensor-Cloud Framework:} Bridging Wireless Sensor Networks (WSNs) with cloud elasticity; virtual sensors, multi-tenant provisioning topologies, and layered middleware architectures.
\end{enumerate}

\section{Containers vs Virtual Machines: Architectural Divergence}
\subsection{Structural Comparison}
Traditional virtualization relies on a hypervisor (Type-1 or Type-2) to emulate physical hardware (CPUs, chipsets, network interface cards, memory controllers, BIOS). Each guest Virtual Machine must run a complete, independent guest operating system including its own kernel, device drivers, and system services.

In contrast, \textbf{Operating System-Level Virtualization (Containerization)} relies on kernel-level isolation features within a single host operating system kernel to partition user space into isolated execution environments known as \textbf{Containers}.

\begin{figure}[H]
\centering
\begin{tikzpicture}[font=\footnotesize]
  % VM Stack
  \node[dark, minimum width=5.5cm, minimum height=0.6cm] (hw1) {Physical Bare-Metal Hardware};
  \node[mid, above=0.2cm of hw1, minimum width=5.5cm, minimum height=0.6cm] (hyp) {Hypervisor (Type-1 / Type-2)};
  
  \node[box, above=0.3cm of hyp.north west, anchor=south west, minimum width=2.5cm, minimum height=0.5cm] (gos1) {Guest OS (Kernel)};
  \node[box, above=0.3cm of hyp.north east, anchor=south east, minimum width=2.5cm, minimum height=0.5cm] (gos2) {Guest OS (Kernel)};
  
  \node[steel, above=0.15cm of gos1, minimum width=2.5cm, minimum height=0.5cm] (bin1) {Bins / Libs};
  \node[steel, above=0.15cm of gos2, minimum width=2.5cm, minimum height=0.5cm] (bin2) {Bins / Libs};
  
  \node[gold, above=0.15cm of bin1, minimum width=2.5cm, minimum height=0.5cm] (app1) {App A};
  \node[gold, above=0.15cm of bin2, minimum width=2.5cm, minimum height=0.5cm] (app2) {App B};

  \node[below=0.1cm of hw1, font=\bfseries] {(a) Virtual Machines (Hardware Virtualization)};

  % Container Stack
  \node[dark, right=2cm of hw1, minimum width=5.5cm, minimum height=0.6cm] (hw2) {Physical Bare-Metal Hardware};
  \node[mid, above=0.2cm of hw2, minimum width=5.5cm, minimum height=0.6cm] (hostos) {Host Operating System (Shared Linux Kernel)};
  \node[dark, above=0.2cm of hostos, minimum width=5.5cm, minimum height=0.6cm] (dockerd) {Docker Engine / Container Runtime};

  \node[steel, above=0.3cm of dockerd.north west, anchor=south west, minimum width=2.5cm, minimum height=0.5cm] (cbin1) {Bins / Libs};
  \node[steel, above=0.3cm of dockerd.north east, anchor=south east, minimum width=2.5cm, minimum height=0.5cm] (cbin2) {Bins / Libs};

  \node[gold, above=0.15cm of cbin1, minimum width=2.5cm, minimum height=0.5cm] (capp1) {Container App A};
  \node[gold, above=0.15cm of cbin2, minimum width=2.5cm, minimum height=0.5cm] (capp2) {Container App B};

  \node[below=0.1cm of hw2, font=\bfseries] {(b) Containers (OS-Level Virtualization)};
\end{tikzpicture}
\caption{Architectural Comparison: Virtual Machines vs Application Containers.}
\end{figure}

\subsection{Quantitative Tradeoff Analysis}
\begin{center}
\begin{tabularx}{\textwidth}{>{\bfseries}l X X}
\toprule
Metric / Dimension & Virtual Machines (VMs) & Containers (Docker)\\
\midrule
Isolation Level & High (Hardware-enforced by MMU, VT-x, and separate kernel ring 0). & Medium-High (Process-level isolation enforced by host kernel namespaces and cgroups).\\
Kernel Instances & Each VM executes its own discrete OS kernel. & All containers share the single host OS kernel.\\
Startup Latency & Minutes ($30\text{ s}$ to $180\text{ s}$) due to full virtual BIOS and OS boot. & Milliseconds to seconds ($< 1\text{ s}$) as processes spawn directly.\\
Disk Footprint & Gigabytes ($2\text{ GB}$ to $40\text{ GB}$) per instance (entire OS image). & Megabytes ($10\text{ MB}$ to $500\text{ MB}$) containing only app dependencies.\\
CPU/RAM Overhead & Significant ($10\%$ to $25\%$) dedicated to running guest OS daemons. & Minimal ($< 1\%$), near bare-metal execution performance.\\
Density & Tens of VMs per physical bare-metal host. & Hundreds to thousands of containers per physical host.\\
Portability & Heavy image formats (VMDK, VHD, QCOW2) tied to hypervisor formats. & Standard OCI image formats portable across any Linux distribution.\\
\bottomrule
\end{tabularx}
\end{center}

\section{The Docker Ecosystem and Architecture}
\subsection{Docker Engine Components}
Docker is an open-source platform that automates the deployment, scaling, and management of applications inside portable container environments. The Docker architecture consists of:
\begin{itemize}[leftmargin=*]
  \item \textbf{Docker Daemon (\texttt{dockerd}):} The persistent background service running on the host that builds images, manages container states, attaches virtual networks, and coordinates storage volumes.
  \item \textbf{Docker Client (\texttt{docker}):} The command-line interface (CLI) tool through which users interact with \texttt{dockerd} via REST API endpoints over UNIX domain sockets (\texttt{/var/run/docker.sock}) or TCP.
  \item \textbf{Docker Images:} Read-only, immutable templates composed of stacked union filesystem layers.
  \item \textbf{Docker Containers:} Runnable instances of a Docker image. When a container starts, a thin, writable copy-on-write (CoW) layer is mounted on top of the immutable image layers.
  \item \textbf{Docker Registry (Docker Hub):} Centralized repositories for publishing, distributing, and pulling container images.
\end{itemize}

\subsection{Copy-on-Write (CoW) and Storage Layering}
Docker utilizes union filesystems (such as \textbf{OverlayFS} or \textbf{AUFS}) to stack directories and present them as a single cohesive unified filesystem:
\begin{itemize}[leftmargin=*]
  \item When multiple containers instantiate the identical base image (\texttt{ubuntu:22.04}), they all share the identical read-only physical disk blocks in memory.
  \item If a container modifies an existing file, the underlying file is copied up from the read-only layer into the container's top-level writable layer before modification (\emph{Copy-on-Write}). This ensures zero mutation of base images.
\end{itemize}

\subsection{Dockerfile Syntax and Construction}
A \textbf{Dockerfile} is a declarative text manifest specifying the sequential instructions required to assemble a Docker image:
\begin{lstlisting}[language=bash]
# Base image layer
FROM python:3.9-slim

# Set execution working directory
WORKDIR /app

# Copy dependency manifest and install packages
COPY requirements.txt .
RUN pip install --no-cache-dir -r requirements.txt

# Copy application source code into image layer
COPY . .

# Expose internal listening port
EXPOSE 5000

# Specify default container entry command
CMD ["python", "app.py"]
\end{lstlisting}

\subsection{Complete Docker Command Lifecycle}
\begin{center}
\begin{tabularx}{\textwidth}{>{\bfseries\ttfamily}l X}
\toprule
Docker Command & Operational Execution Description\\
\midrule
docker build -t myapp:1.0 . & Reads the local Dockerfile, executes instructions sequentially, caches intermediate layers, and tags the resulting image.\\
docker run -d -p 8080:5000 myapp:1.0 & Instantiates a container in detached mode (\texttt{-d}), binding host port 8080 to container port 5000 (\texttt{-p}).\\
docker ps -a & Lists all containers (both actively running and terminated/exited).\\
docker stop <container\_id> & Sends SIGTERM followed by SIGKILL to gracefully terminate container processes.\\
docker rm <container\_id> & Deletes a stopped container instance and purges its writable layer.\\
docker rmi <image\_id> & Permanently deletes an unreferenced container image from the local storage cache.\\
docker network ls & Displays all configured virtual bridge, overlay, and host networks.\\
docker exec -it <id> /bin/bash & Spawns an interactive bash shell session inside an actively running container.\\
\bottomrule
\end{tabularx}
\end{center}

\section{Container Orchestration: Kubernetes and etcd}
When applications scale to thousands of containers spanning hundreds of bare-metal servers, manual container management becomes impossible. \textbf{Kubernetes} automates container scheduling, automated rollouts, service discovery, load balancing, and self-healing.

\begin{defbox}{Role of \texttt{etcd} in Kubernetes Clusters}
\textbf{\texttt{etcd}} is a strongly consistent, distributed, and highly available key-value datastore based on the \textbf{Raft consensus algorithm}.\\
In Kubernetes, \texttt{etcd} serves as the ultimate single source of truth for the entire cluster. It persistently stores:
\begin{itemize}[leftmargin=*]
  \item Cluster configuration states, namespaces, and node registrations.
  \item Pod definitions, ReplicaSets, and Deployments.
  \item Secrets, ConfigMaps, and cluster metadata.
\end{itemize}
The Kubernetes control plane components (API Server, Scheduler, Controller Manager) are stateless; they coordinate state transitions strictly by reading from and writing to \texttt{etcd}.
\end{defbox}

\section{Sensor-Cloud Infrastructure}
\subsection{Limitations of Traditional Wireless Sensor Networks (WSNs)}
Traditional Wireless Sensor Networks (WSNs) consist of small sensor motes deployed to monitor physical environments (temperature, moisture, seismic activity). However, conventional WSNs exhibit severe structural limitations:
\begin{itemize}[leftmargin=*]
  \item \textbf{Isolated Application Silos:} A sensor network deployed for soil moisture monitoring is typically accessible only to a single proprietary backend application.
  \item \textbf{Resource Exhaustion:} Sensor nodes possess limited battery life, micro-controllers with few kilobytes of RAM, and low-power radios.
  \item \textbf{No Multi-Tenancy:} Multiple users with differing analytical queries cannot dynamically share the same physical sensor hardware simultaneously.
\end{itemize}

\subsection{The Sensor-Cloud Paradigm}
\begin{defbox}{Sensor-Cloud Definition}
\textbf{Sensor-Cloud} is an infrastructure that virtualizes physical Wireless Sensor Networks into dynamically manageable, reconfigurable, and multi-tenant \textbf{Virtual Sensors} provisioned on-demand over the cloud.
\end{defbox}

\begin{figure}[H]
\centering
\begin{tikzpicture}[node distance=1.2cm, font=\small]
  % Layers
  \node[dark, minimum width=10cm, minimum height=0.9cm] (app) {Application Layer (End Users, Smart Agriculture, Healthcare, Telematics)};
  \node[mid, below=0.5cm of app, minimum width=10cm, minimum height=0.9cm] (scsp) {Sensor-Cloud Middleware / Virtualization Layer (SCSP, Virtual Sensor Provisioning)};
  \node[gold, below=0.5cm of scsp, minimum width=10cm, minimum height=0.9cm] (gw) {Gateway / Proxy Layer (Protocol Normalization, Data Aggregation)};
  \node[steel, below=0.5cm of gw, minimum width=10cm, minimum height=0.9cm] (wsn) {Physical WSN Layer (Heterogeneous Sensors, Motes, ZigBee, LoRaWAN)};

  \draw[arr, <->] (app) -- (scsp);
  \draw[arr, <->] (scsp) -- (gw);
  \draw[arr, <->] (gw) -- (wsn);
\end{tikzpicture}
\caption{Layered Architecture of Sensor-Cloud Computing.}
\end{figure}

\subsection{Stakeholders and Economic Model}
Sensor-Cloud establishes a three-party multi-tenant marketplace:
\begin{enumerate}[leftmargin=*]
  \item \textbf{Physical Sensor Owners:} Entities that procure and deploy physical sensor hardware in the field. They register their WSNs with the Sensor-Cloud provider to monetize idle sensing capacity.
  \item \textbf{Sensor-Cloud Service Provider (SCSP):} The cloud entity that maintains the virtualization middleware, registers physical sensors, provisions virtual sensors, negotiates SLAs, and handles metering and billing.
  \item \textbf{End Users:} Consumers or enterprises who provision virtual sensors on-demand for specific temporal durations and geographic boundaries without purchasing physical hardware.
\end{enumerate}

\subsection{Virtual Sensor Configurations and Mapping Topologies}
A \textbf{Virtual Sensor} is a software abstraction representing a sensing unit provisioned in the cloud. It decouples the user application from the underlying physical sensor hardware. The virtualization engine supports three primary configurations:

\begin{itemize}[leftmargin=*]
  \item \textbf{One-to-One Mapping ($1:1$):} A single virtual sensor directly represents a single physical sensor mote.
  \item \textbf{Many-to-One Mapping ($N:1$):} Multiple physical sensor motes are aggregated into a single virtual sensor. For example, averaging temperature readings from 20 distributed agricultural sensor nodes across a field into a single virtual temperature feed.
  \item \textbf{One-to-Many Mapping ($1:N$):} A single high-capability physical sensor is dynamically partitioned into multiple virtual sensors serving different end users simultaneously with distinct sampling frequencies and delivery thresholds.
  \item \textbf{Spatio-Temporal Virtual Sensors:} Virtual sensors synthesized dynamically to cover an arbitrary geographical polygon $\mathcal{R}$ over a designated time window $[t_1, t_2]$, collecting and synthesizing telemetry from all physical sensors falling inside $\mathcal{R}$.
\end{itemize}

\section{Comprehensive Solved Assignment Problems (Week 8)}

\begin{solved}{Docker Image Deletion Command (Assignment 8, Q1)}
\Given Which command deletes a Docker image that is no longer required?
\Steps
\begin{enumerate}[leftmargin=*]
  \item \texttt{docker rm} removes stopped \emph{containers}.
  \item \texttt{docker rmi} (\textbf{r}e\textbf{m}ove \textbf{i}mage) purges container \emph{images}.
\end{enumerate}
\Answer \textbf{Option B: \texttt{docker rmi}}.
\end{solved}

\begin{solved}{Listing Docker Virtual Networks (Assignment 8, Q2)}
\Given Identify the command used to display all available Docker networks on the host.
\Steps
\begin{enumerate}[leftmargin=*]
  \item Docker CLI manages network namespaces via the \texttt{network} command group.
  \item The command \texttt{docker network ls} lists all bridge, host, and overlay networks.
\end{enumerate}
\Answer \textbf{Option A: \texttt{docker network ls}}.
\end{solved}

\begin{solved}{VM vs Container Kernel Isolation Comparison (Assignment 8, Q3)}
\Given Which is the most accurate comparison between VMs and containers?
\Steps
\begin{enumerate}[leftmargin=*]
  \item Virtual Machines virtualize physical hardware, executing complete guest operating systems with independent kernels.
  \item Containers share the underlying host OS kernel, providing process isolation via Linux kernel namespaces and cgroups.
\end{enumerate}
\Answer \textbf{Option B: Containers generally share the host kernel, while VMs virtualize hardware and run separate guest OS kernels}.
\end{solved}

\begin{solved}{Primary Function of etcd in Kubernetes (Assignment 8, Q4)}
\Given What is the fundamental responsibility of \texttt{etcd} in a Kubernetes cluster?
\Steps
\begin{enumerate}[leftmargin=*]
  \item \texttt{etcd} is a distributed key-value store used to persist cluster configuration data, metadata, and state information across all master control plane components.
\end{enumerate}
\Answer \textbf{Option B: Storing Kubernetes cluster state and configuration data}.
\end{solved}

\begin{solved}{Hypervisor Type Classification (Assignment 8, Q5)}
\Given Identify an example of a Type 2 (hosted) hypervisor among: VMware ESXi, Xen, Oracle VirtualBox, Microsoft Hyper-V Server.
\Steps
\begin{enumerate}[leftmargin=*]
  \item VMware ESXi, Xen, and Hyper-V Server are Type-1 bare-metal hypervisors.
  \item \textbf{Oracle VirtualBox} runs as an application process on top of a standard host operating system (Windows, Linux, macOS), making it a Type 2 hypervisor.
\end{enumerate}
\Answer \textbf{Option C: Oracle VirtualBox}.
\end{solved}

\begin{solved}{Energy Reduction via VM Consolidation (Assignment 8, Q6)}
\Given In a data center of 100 servers, 40 are heavily loaded while 60 operate at very low utilization. What is the most effective energy saving strategy?
\Steps
\begin{enumerate}[leftmargin=*]
  \item Idle servers consume $50\%$ to $60\%$ of peak power.
  \item Consolidating VMs from underutilized servers onto fewer hosts via live migration allows the empty physical machines to be powered down into sleep states, maximizing energy savings.
\end{enumerate}
\Answer \textbf{Option B: Consolidate VMs onto fewer servers and switch idle servers to a low-power state}.
\end{solved}

\begin{solved}{Role of Middleware in Sensor-Cloud (Assignment 8, Q7)}
\Given Why is a middleware layer required between Wireless Sensor Networks (WSNs) and cloud applications?
\Steps
\begin{enumerate}[leftmargin=*]
  \item Physical sensor motes use diverse, proprietary communication protocols and data formats.
  \item The middleware abstracts heterogeneous physical sensors, enabling cloud-accessible virtual sensors with standardized APIs.
\end{enumerate}
\Answer \textbf{Option B: To abstract heterogeneous sensors and provide cloud-accessible sensor resources}.
\end{solved}

\begin{solved}{Virtual Sensor Provisioning Topologies (Assignment 8, Q8)}
\Given Can multiple physical sensors be aggregated to form a single virtual sensor?
\Steps
\begin{enumerate}[leftmargin=*]
  \item Yes, Sensor-Cloud supports Many-to-One ($N:1$) mapping, allowing data from multiple geographic sensors to be aggregated into a single virtual sensor feed.
\end{enumerate}
\Answer \textbf{Option A: Yes, Many-to-One mapping is supported}.
\end{solved}

\begin{solved}{Docker Container Isolation Primitives (Assignment 8, Q9)}
\Given Which Linux kernel primitives provide resource metering and process isolation for Docker containers?
\Steps
\begin{enumerate}[leftmargin=*]
  \item \textbf{Namespaces} partition system resources (process IDs, network interfaces, mount points).
  \item \textbf{Control Groups (cgroups)} meter and limit hardware consumption (CPU shares, memory quotas, I/O bandwidth).
\end{enumerate}
\Answer \textbf{Option C: Namespaces and cgroups}.
\end{solved}

\begin{solved}{Economic Benefit of Sensor-Cloud to Sensor Owners (Assignment 8, Q10)}
\Given How do physical sensor owners benefit economically in the Sensor-Cloud model?
\Steps
\begin{enumerate}[leftmargin=*]
  \item Sensor owners register their physical WSNs with the Sensor-Cloud Service Provider (SCSP) and earn recurring revenue from multi-tenant subscriptions without building custom software applications.
\end{enumerate}
\Answer \textbf{Option B: Monetizing deployed physical sensor capacity through multi-tenant usage fees}.
\end{solved}
